\chapter{Eliminación física}
\label{cap:eliminacion}

\begin{importante}
Lo que se elimina aquí \textbf{se borra para siempre} junto con todo lo que depende de ello, y no se puede recuperar. Solo el DIRECTORIO entra a esta pantalla, y ningún cargo puede recibir este permiso.
\end{importante}

\section{Qué se puede eliminar}

\captura{eliminacion}{Eliminación física}

Solo se puede elegir lo que \textbf{ya está dado de baja} y las \textbf{ventas ya cobradas}. Hay una pestaña por tipo: empleados, cargos, secciones y mesas, productos, categorías, insumos, promociones y ventas. Cada pestaña dice cuántos registros tiene marcados.

\begin{enumerate}
  \item Entre a \pantalla{Administración → Eliminación física}.
  \item En cada pestaña, toque las filas que quiere borrar. \boton{Marcar todos} y \boton{Quitar marcas} ayudan con muchas filas. En \textbf{Ventas} primero se eligen las fechas.
  \item Se puede marcar en varias pestañas y borrar todo junto con \boton{Eliminar N marcados}.
\end{enumerate}

\section{Qué arrastra cada cosa (en cascada)}

Al eliminar algo se borra también todo lo que depende de eso. Cada fila avisa lo que arrastra, por ejemplo \textit{2 ventas se borran también}.

\begin{table}[H]
  \centering
  \caption{Lo que se borra junto con cada registro}
  \rowcolors{2}{hueso}{white}
  \begin{tabularx}{\linewidth}{lX}
    \toprule
    \textbf{Si elimina…} & \textbf{También se borra} \\
    \midrule
    Un cargo & Sus empleados \\
    Un empleado & Las ventas que abrió, a las que agregó un envío o que cobró \\
    Una sección & Sus mesas y las ventas de esas mesas \\
    Una categoría & Sus subcategorías y sus productos \\
    Un producto & Su receta, las promociones que lo incluyen y las ventas donde aparece \\
    Una promoción & Las ventas donde se vendió \\
    Un insumo & Su lugar en las recetas, sus movimientos y las ventas donde se quitó \\
    Una venta & Sus productos, envíos, pagos e impresiones del ticket \\
    \bottomrule
  \end{tabularx}
\end{table}

\section{Confirmar}

\captura[0.8]{eliminacion-confirmar}{Confirmación de la eliminación}

Antes de borrar, el sistema muestra \textbf{todo lo que se va}, incluido lo que arrastra. Para seguir hay que escribir la palabra \texttt{ELIMINAR} y presionar \boton{Sí, eliminar}.

Si algo de lo elegido sigue activo, o si arrastra una mesa con un pedido abierto, el sistema dice qué lo impide y no deja seguir: dé de baja ese registro o cobre el pedido antes.

\begin{consejo}
Antes de una eliminación grande, descargue el PDF de Reportes del período que se va a borrar: las ventas eliminadas dejan de aparecer en los reportes.
\end{consejo}
